\documentclass[12pt]{article}

\usepackage{sbc-template} \usepackage{graphicx} \usepackage{url} \usepackage[brazil]{babel}
\usepackage[utf8]{inputenc}

\sloppy

\title{InstanciaIris: Pipeline de Machine Learning com Rastreamento de Experimentos em MLflow e API de Predição}

\author{Diego Cordeiro de Oliveira\inst{1}}

\address{Departamento de Computação -- Universidade Federal do Piauí (UFPI)\\  Teresina -- PI -- Brasil \email{diego.cordeiro@ifpi.edu.br} }

\begin{document}

\maketitle

\begin{abstract}
This technical report presents the instantiation and the end-to-end execution of the InstanciaIris project,
a machine learning workflow built on the FabricaIA template. The goal was to establish a reproducible
pipeline for the Iris dataset, covering preprocessing (label encoding and feature standardization),
training of a Random Forest classifier with experiment tracking in MLflow, model persistence, and the
exposure of a REST prediction endpoint via FastAPI. The notebook was executed without errors: the
Random Forest model (100 estimators, maximum depth 10) achieved an accuracy of 0.9333 and an F1 score
of 0.90, being registered in the MLflow experiment testeprojeto\_experiments. The prediction API
answered a test request returning a class and the associated probabilities. A known limitation is
highlighted: the /predict endpoint does not reapply the feature scaler; as future work, persisting the
preprocessing pipeline and applying it during inference is recommended.
\end{abstract}

\begin{resumo}
Este relatório técnico apresenta a instanciação e a execução ponta a ponta do projeto InstanciaIris,
um fluxo de aprendizado de máquina construído sobre o template FabricaIA. O objetivo foi estabelecer
um pipeline reproduzível para o dataset Iris, abrangendo pré-processamento (codificação do rótulo e
padronização das características), treinamento de um classificador Random Forest com rastreamento de
experimentos em MLflow, persistência do modelo e exposição de um endpoint de predição REST via FastAPI.
O notebook foi executado sem erros: o modelo Random Forest (100 estimadores, profundidade máxima 10)
alcançou acurácia de 0,9333 e F1 de 0,90, sendo registrado no experimento testeprojeto\_experiments do
MLflow. A API de predição respondeu a uma requisição de teste com a classe e as probabilidades
associadas. Destaca-se a limitação conhecida: o endpoint /predict não reaplica o escalador de
características; como trabalho futuro, sugere-se persistir o pipeline de pré-processamento e aplicá-lo
na inferência.
\end{resumo}

\section{Introdução}
Projetos de aprendizado de máquina não se encerram no treinamento de um modelo: a reprodutibilidade,
o rastreamento de experimentos e a disponibilização do modelo para inferência são etapas igualmente
relevantes para consolidar uma solução confiável. Nesse contexto, o uso de templates e de ferramentas de
experimentação contribui para padronizar o fluxo de trabalho e reduzir o risco de perda de informação
sobre versões, hiperparâmetros e métricas.

O dataset Iris, proposto por Fisher, é um benchmark clássico de classificação supervisionada, com 150
amostras e três classes de espécies de flor, sendo adequado para validar pipelines de pré-processamento,
treinamento e avaliação.

Diante disso, este relatório documenta a instanciação e a execução do projeto InstanciaIris, com os
seguintes objetivos: (i) instanciar um projeto de aprendizado de máquina a partir do template
FabricaIA; (ii) construir um pipeline de pré-processamento e treinamento com rastreamento de
experimentos em MLflow; (iii) expor um modelo treinado por meio de uma API REST com FastAPI; e
(iv) validar o funcionamento ponta a ponta por meio de testes automatizados e de uma chamada de
predição real.

\section{Fundamentação Teórica}
O aprendizado de máquina supervisionado busca aprender uma função que mapeia características de entrada
para um rótulo de saída a partir de exemplos rotulados. O Random Forest é um método de conjunto (ensemble)
baseado em árvores de decisão, no qual múltiplas árvores são construídas sobre subamostras aleatórias dos
dados e de suas características, e a predição final é obtida por votação majoritária~\cite{breiman2001randomforests, dietterich2000ensemble}.
A implementação utilizada neste trabalho é a fornecida pela biblioteca scikit-learn~\cite{pedregosa2011sklearn}.

O dataset Iris~\cite{fisher1936iris} contém 150 amostras organizadas em três classes (setosa, versicolor
e virginica), caracterizadas por quatro medidas morfológicas: comprimento e largura da sépala e
comprimento e largura da pétala.

Em tarefas de classificação, o pré-processamento é tipicamente composto pela codificação de variáveis
categóricas e pela padronização de características numéricas. A codificação de rótulos (LabelEncoder)
transforma as classes em valores inteiros, enquanto a padronização (StandardScaler)~\cite{pedregosa2011sklearn} ajusta cada
característica para média zero e desvio padrão unitário, evitando que a diferença de escala influencie
o modelo.

O rastreamento de experimentos é fundamental para a reprodutibilidade em aprendizado de máquina. O MLflow~
\cite{zaharia2018mlflow} permite registrar experimentos, parâmetros, métricas e artefatos de modelos,
organizados em execuções (runs) agrupadas em experimentos. Neste trabalho, o MLflow foi configurado com
backend SQLite local (sqlite:///mlflow.db).

Para a avaliação de modelos de classificação, métricas como acurácia, precisão, recall e F1-score, derivadas da matriz de confusão, são amplamente utilizadas~\cite{sokolova2009systematic}. A divisão dos dados em conjuntos de treino e teste é uma prática padrão de validação, e a estratificação é comumente empregada para preservar a distribuição das classes~\cite{kohavi1995cross}. Na busca por melhor desempenho, a otimização de hiperparâmetros, por busca aleatória ou em grade, é frequentemente utilizada~\cite{bergstra2012random}.

Para disponibilizar o modelo, utiliza-se uma API REST construída com FastAPI~
\cite{ramirezfastapi}, framework web assíncrono que emprega modelos de dados definidos com Pydantic para validar as requisições.

\section{Metodologia}

\subsection{Ambiente e estrutura do projeto}
O projeto foi instanciado em um ambiente virtual Python 3.12.9, com as bibliotecas scikit-learn 1.9.0,
MLflow 3.16.0, pandas, numpy e FastAPI. A estrutura de diretórios segue o template FabricaIA, com as
pastas config/, data/, models/, notebooks/, src/, tests/ e docs/.

\subsection{Dados e pré-processamento}
O arquivo data/raw/iris.csv contém 150 linhas e 5 colunas (quatro características numéricas e a coluna
alvo target). O pipeline de pré-processamento envolve limpeza, codificação do rótulo e padronização das
características, seguido da divisão dos dados em treino e teste (80/20), resultando em 119 amostras de
treino.

\begin{verbatim}
df = load_data(project_path("data","raw","iris.csv"))
df_clean = clean_data(df, drop_duplicates=True)
df_encoded = encode_categorical(df_clean)   # alvo -> 0/1/2
df_scaled = scale_features(df_encoded)      # 4 features (StandardScaler)
X_train, X_test, y_train, y_test = split_data(df_scaled, "target")
\end{verbatim}

\subsection{Treinamento e rastreamento de experimentos}
O classificador Random Forest foi configurado com 100 estimadores e profundidade máxima 10, treinado por
meio do componente ModelTrainer do projeto. O rastreamento foi realizado com MLflow, no experimento
testeprojeto\_experiments, com backend SQLite local, registrando parâmetros (n\_estimators, max\_depth e
número de amostras), métricas e o modelo como artefato.

\begin{verbatim}
model = trainer.train_model(X_train, y_train, "random_forest",
    run_name="random_forest_baseline", n_estimators=100, max_depth=10)
metrics = trainer.evaluate_model(model, X_test, y_test)
trainer.save_model(model, project_path("models","trained","model.pkl"),
    artifact_path="random_forest_model")
trainer.end_run()
\end{verbatim}

\subsection{API de predição}
Na inicialização, a aplicação FastAPI carrega os modelos de models/trained/*.pkl e deriva o nome do
modelo removendo o sufixo \_model; como o arquivo persistido é model.pkl, o nome válido é model. O
endpoint POST /predict recebe um corpo JSON com as características e o nome do modelo e retorna a classe
predita, as probabilidades e a confiança.

\begin{verbatim}
POST /predict
{ "features": { "sepal_length": 5.1, "sepal_width": 3.5,
    "petal_length": 1.4, "petal_width": 0.2 },
  "model_name": "model" }
\end{verbatim}

A resposta retornada é exibida abaixo.

\begin{verbatim}
{ "prediction": 2,
  "probability": { "0": 0.01, "1": 0.24, "2": 0.75 },
  "confidence": 0.75 }
\end{verbatim}

\subsection{Validação por testes automatizados}
A suíte de testes do projeto foi executada e validou os componentes de processamento, treinamento,
engenharia de características e visualização, com 18 testes aprovados. Como o pacote não está instalado
de forma editável no ambiente, os testes foram executados com PYTHONPATH apontando para a raiz do projeto.

\section{Resultados e Discussão}
Os resultados obtidos confirmam a execução do pipeline de ponta a ponta. A Tabela~\ref{tab:metricas}
resume as métricas do modelo Random Forest treinado, registradas no experimento testeprojeto\_experiments.

\begin{table}[h]
\centering
\caption{Métricas registradas no MLflow (precisão, recall e F1 para a classe 1).}
\label{tab:metricas}
\begin{tabular}{lcccc}
\hline
Run & Acurácia & Precisão & Recall & F1 \\
\hline
random\_forest\_baseline & 0{,}9333 & 0{,}90 & 0{,}90 & 0{,}90 \\
entrega\_topicos\_avancados & 0{,}9333 & 0{,}90 & 0{,}90 & 0{,}90 \\
\hline
\end{tabular}
\end{table}

O modelo alcançou acurácia de 93{,}33\% e F1 de 0{,}90, resultado coerente com a simplicidade e a
separabilidade do dataset Iris. Como esperado para um modelo baseline, não foi realizada otimização de
hiperparâmetros; a escolha de 100 estimadores e profundidade máxima 10 forneceu um desempenho estável.
O modelo foi persistido em models/trained/model.pkl e registrado como artefato no experimento
testeprojeto\_experiments.

\begin{table}[h]
\centering
\caption{Predições com características padronizadas (0=setosa, 1=versicolor, 2=virginica).}
\label{tab:predicoes}
\begin{tabular}{lccc}
\hline
Espécie & Características cruas & Características padronizadas & Predição \\
\hline
setosa & 5{,}1, 3{,}5, 1{,}4, 0{,}2 & -1{,}024, -0{,}850, -1{,}368, -1{,}296 & 0 \\
versicolor & 7{,}0, 3{,}2, 4{,}7, 1{,}4 & 0{,}288, -0{,}873, 0{,}228, 0{,}355 & 1 \\
virginica & 6{,}3, 3{,}3, 6{,}0, 2{,}5 & -0{,}234, -0{,}709, 1{,}013, 1{,}032 & 2 \\
\hline
\end{tabular}
\end{table}

A suíte de testes passou integralmente (18 testes aprovados), corroborando a integridade dos componentes.

\section{Conclusão}
Este relatório documentou a instanciação e a execução do projeto InstanciaIris. Os objetivos foram
alcançados: o projeto foi instanciado a partir do template FabricaIA; um pipeline de pré-processamento e
treinamento foi construído e executado com sucesso no notebook; o modelo Random Forest foi treinado,
avaliado e rastreado em MLflow, com acurácia de 93{,}33\% e F1 de 0{,}90; o modelo foi persistido e exposto
por uma API REST via FastAPI; e o funcionamento foi validado por testes automatizados e por uma chamada
de predição real.

Entre as contribuições, destacam-se a definição de um fluxo de trabalho reproduzível, o uso de MLflow
para a rastreabilidade de experimentos e a disponibilização de um serviço de inferência.

\bibliographystyle{sbc}
\bibliography{sbc-template}

\end{document}
